% TOP_PROBLEM: {"problemId": "problem.kaplansky-s-direct-finiteness-conjecture", "canonicalTitle": "Kaplansky's Direct Finiteness Conjecture", "releaseRank": 440, "primaryDomain": "algebra_representation_category", "primaryDomainLabel": "Algebra, representation & category theory", "displayStatus": "open_with_solved_subcases", "releaseStatus": "open_with_solved_subcases"}
% !TeX program = lualatex
% ATTEMPT_STATUS: unresolved
% Research continuation by GPT_6_Astra_Ultra, 27 September 2026.
\documentclass[11pt]{article}
\usepackage[margin=25mm]{geometry}
\usepackage{fontspec}
\setmainfont{Latin Modern Roman}
\usepackage{amsmath,amssymb,mathtools}
\usepackage{unicode-math}
\setmathfont{Latin Modern Math}
\usepackage{xurl,hyperref}
\hypersetup{colorlinks=true,urlcolor=blue}
\setcounter{secnumdepth}{0}
\setlength{\parindent}{0pt}
\setlength{\parskip}{5pt}
\setlength{\emergencystretch}{3em}
\begin{document}
\section*{\texorpdfstring{Kaplansky's Direct Finiteness Conjecture}{Kaplansky's Direct Finiteness Conjecture}}\label{440}
\subsection{Definitions and mathematical statement}
For a group $G$ and field $K$, the group ring $K[G]$ consists of finite sums $\sum_g a_g g$, with coefficientwise addition and multiplication extended bilinearly from the group law. Its identity is $1_K e_G$. A unital ring is directly finite if a right inverse is always a left inverse. The assertion is
\[
 \forall G\ \forall K\ \forall a,b\in K[G],
 \qquad ab=1\Longrightarrow ba=1.
\]
The group may have torsion and the field may have any characteristic. Stable finiteness means direct finiteness of every matrix ring $M_n(K[G])$. The catalogue records a known equivalence of the universal direct- and stable-finiteness targets; the displayed problem itself is the direct-finiteness assertion, with no bound on the supports of $a,b$.
\par\smallskip\textit{Formulation and concept references:} \href{https://web.ma.utexas.edu/users/juschenko/files/soficgroups.pdf}{[S1]}.

\subsection{Short English statement}
In every group ring over a field, must reversing a product that equals one still give one?
\subsection{Research attempt}
\textit{GPT-6 Astra Ultra; 27 September 2026. Status: UNRESOLVED.}

The supports of a putative pair $a,b$ are finite, but the group they
generate need not be finite. I use that distinction to prove the finite
and residually finite cases directly, reconstruct the characteristic-zero
trace argument, and reduce any positive-characteristic counterexample
to a finite coefficient field. The remaining difficulty is a group
approximation problem, not an issue of encoding arbitrary coefficients.
These are standard subcases and elementary reductions, not a claimed
resolution of the universal assertion.

\paragraph{Finite groups and exact finite quotients.}
If $G$ is finite, left multiplication represents $K[G]$ on the
$|G|$-dimensional vector space $K[G]$. The identity $ab=1$ gives
$L_aL_b=I$, and finite-dimensional linear algebra gives $L_bL_a=I$.
Applying the latter to $1$ gives $ba=1$. No characteristic restriction,
semisimplicity, or division by $|G|$ is involved. Thus torsion and the
case $\operatorname{char}K\mid|G|$ do not invalidate this argument.

Suppose now that $G$ is residually finite and $ab=1$. Let $S$ be the
finite support of $ba-1$. For every pair of distinct elements of $S$,
residual finiteness gives a finite quotient separating them. Taking
the product of these finitely many quotients gives one finite quotient
$q:G\to Q$ injective on $S$. The induced group-ring map sends $ab$ to
$1$, so the finite-group result sends $ba-1$ to zero. But its distinct
support elements have distinct images, and their coefficients cannot
cancel. Hence every coefficient is zero already and $ba=1$.

More generally, only the subgroup generated by the supports of $a,b$
matters. All products are computed in that subgroup's group ring. Thus
the universal problem reduces to finitely generated groups, but finite
generation supplies no separating finite quotients. Replacing
``finitely generated'' by ``finite'' would discard exactly the hard
case of this argument.

\paragraph{The trace proof in characteristic zero.}
For $K=\mathbb C$, view the group ring as bounded left-convolution
operators on $\ell^2(G)$ and place it in the group von Neumann algebra.
This algebra has the finite trace
\[
 \tau(T)=\langle T\delta_e,\delta_e\rangle,
 \qquad\tau(I)=1.
\]
It is faithful: if $\tau(T^*T)=0$, then $T\delta_e=0$; commuting with
right translations gives $T\delta_g=0$ for every $g$, hence $T=0$.
The tracial identity $\tau(UV)=\tau(VU)$ extends the elementary
coefficient-of-identity identity from convolution operators. The
von Neumann algebra and polar-decomposition facts are the standard
analytic inputs to this known proof, discussed in [S1].

From $ab=I$ one has
\[
 I=abb^*a^*\le\|b\|^2aa^*.
\]
Thus $aa^*$ is invertible. In the polar decomposition $a=u|a|$ the
final support of $u$ is $uu^*=I$. The projection $u^*u$ has the same
trace as $uu^*$, so faithfulness gives $u^*u=I$ as well. The initial
support being $I$ means that $a$ has zero kernel. Since
$a(ba-I)=0$, this injectivity implies $ba=I$. This is the step that
fails in an arbitrary infinite-dimensional operator algebra: its
projections need not be controlled by a faithful finite trace.

For an arbitrary characteristic-zero field $K$, the finitely many
coefficients in $a,b$ generate a finitely generated extension $K_0$ of
$\mathbb Q$. It embeds into $\mathbb C$: choose algebraically
independent complex images of a transcendence basis and extend across
the remaining finite algebraic extension. The induced group-ring map
is injective. The complex result therefore proves $ba=1$ over $K_0$
and then over $K$. This reduction uses only the coefficient field of
the particular pair, not an embedding of every characteristic-zero
field into $\mathbb C$.

\paragraph{Positive characteristic: finite coefficient fields suffice.}
Assume there is a counterexample over a field $K$ of characteristic
$p$. List the finite supports of $a$ and $b$ and replace their
coefficients by variables. The condition $ab=1$ is a finite system of
polynomial equations over $\mathbb F_p$, one equation for each group
element appearing in the finite product support or at the identity.
Choose a coefficient $c$ of $ba-1$ which is nonzero in the given
counterexample, and add a variable $z$ and the equation
\[
 zc-1=0.                                               \tag{1}
\]
The resulting ideal is proper because the original coefficients and
$z=c^{-1}$ give a point over $K$. A maximal ideal above it has residue
field a finite algebraic extension of $\mathbb F_p$, by Zariski's
lemma. Specializing to that finite field preserves all the equations
$ab=1$, while (1) preserves a nonzero coefficient of $ba-1$.
We obtain a counterexample over a finite field, in the same group.

Some coefficients may specialize to zero, so the supports may shrink;
this causes no problem. The chosen witness coefficient remains nonzero.
The conclusion is a reduction to finite \emph{fields}, not finite
\emph{groups}. The latter are already settled and cannot be obtained
without additional information about the group law on these supports.
An effective procedure for finding a quotient is also not supplied by
this purely existential specialization argument.

\paragraph{A countertest to bare operator injectivity arguments.}
On a vector space with basis $e_0,e_1,\ldots$, let $S e_j=e_{j+1}$
and let $T e_0=0$, $T e_{j+1}=e_j$. Then
\[
 TS=I,\qquad ST=I-P_0\ne I.
\]
Both operators have extremely simple formulas, but the endomorphism
ring is not directly finite. This is not a group-ring counterexample:
no representation as finite convolution sums in a common group ring
has been given. It demonstrates exactly why finite support in a
description of an operator is not a substitute for the finite trace
or finite-quotient argument above.

\paragraph{Verification and first gap.}
The finite-dimensional proof was checked without assuming Maschke's
theorem, and the specialization retains an explicit nonzero witness.
The unilateral-shift identities are exact on every basis vector;
truncating them to a finite matrix would introduce a boundary defect
and destroy $TS=I$. [S1] proves a broader positive result for sofic
groups; this attempt does not prove that an arbitrary group has that
approximation property. After the reductions, the first gap is to
exclude a counterexample over a finite field in a finitely generated
group without a known suitable finite approximation. No such exclusion
or counterexample is established. The universal conjecture remains
\textbf{UNRESOLVED}.

\subsection{Sources}
\begin{itemize}
\item[S1]Kate Juschenko, Sofic Groups, Chapter 5: Some conjectures that are valid for sofic groups.
\url{https://web.ma.utexas.edu/users/juschenko/files/soficgroups.pdf}.
\textit{Formulation source.}
\item[E1]Kate Juschenko, \emph{Sofic Groups}, Chapter 5, Section 1.
\url{https://web.ma.utexas.edu/users/juschenko/files/soficgroups.pdf}.
\textit{Primary author exposition corresponding to [S1], consulted for
the characteristic-zero and sofic-group scope on 27 September 2026.
Finite-field specialization is explicitly justified in this attempt.}

\end{itemize}
\end{document}
